\documentclass[11pt, a4paper]{article}

% ── Packages ──────────────────────────────────────────────────────────
\usepackage[utf8]{inputenc}
\usepackage[T1]{fontenc}
\usepackage{lmodern}
\usepackage[margin=1in]{geometry}
\usepackage{amsmath, amssymb, amsthm}
\usepackage{booktabs}
\usepackage{longtable}
\usepackage{graphicx}
\usepackage{hyperref}
\usepackage{xcolor}
\usepackage{listings}
\usepackage{tcolorbox}
\usepackage{polya}

\hypersetup{
  colorlinks=true,
  linkcolor=polyablue,
  urlcolor=polyablue,
  citecolor=polyablue,
}

% ── Code listing style ────────────────────────────────────────────────
\lstdefinestyle{polya-python}{
  language=Python,
  basicstyle=\ttfamily\small,
  keywordstyle=\color{polyablue}\bfseries,
  commentstyle=\color{polyagray}\itshape,
  stringstyle=\color{polyagreen},
  numbers=left,
  numberstyle=\tiny\color{polyagray},
  numbersep=8pt,
  frame=single,
  framerule=0.4pt,
  rulecolor=\color{polyagray},
  breaklines=true,
  showstringspaces=false,
  tabsize=4,
}
\lstset{style=polya-python}
\title{Pareto Optimization -- Product Design Trade-offs}
\author{uber-polya}
\date{2026-02-26}

\begin{document}

\maketitle
\tableofcontents
\newpage

% ── Phase A: Problem Formulation ─────────────────────────────────────
\section{Problem Formulation}

\begin{polyamodel}

\textbf{Problem Type}: Find \hfill
\textbf{Domain}: Multi-Objective Optimization


\end{polyamodel}

% ── Universe ──────────────────────────────────────────────────────────
\subsection{Universe}

\begin{itemize}
  \item $Knee Point: \{cost, inv_durability, thickness, alloy_grade, coating_level, frontier_index\}$
  \item $Utopia: \{cost, inv_durability\}$
  \item $Nadir: \{cost, inv_durability\}$
\end{itemize}

% ── Variables ─────────────────────────────────────────────────────────
\subsection{Variables}

\begin{longtable}{lll}
\toprule
\textbf{Symbol} & \textbf{Meaning} & \textbf{Type / Range} \\
\midrule
\endhead
$x$ & decision variable & $\mathbb{{R}}$ \\
\bottomrule
\end{longtable}

% ── Structure ─────────────────────────────────────────────────────────
\subsection{Structure}

A manufacturing company designs a product with two competing objectives: minimize cost and maximize durability (equivalently, minimize inverse durability). There are 3 design variables -- material thickness, alloy grade, and coating level -- each with physical bounds. The goal is to find the Pareto frontier (the set of non-dominated designs), identify the knee point (best compromise), and quantify the trade-off rates along the frontier.

**Design Variables**:
- Material thickness: [1.0, 10.0] mm
- Alloy grade: [1.0, 5.0] (index)
- Coating level: [0.5, 3.0] (index)

**Objectives**:
- Minimize cost: f1(x) = 2*thickness + 5*alloy\textasciicircum{}2 + 3*coating
- Minimize inverse durability: f2(x) = -ln(thickness) - 2*alloy - 1.5*coating\textasciicircum{}0.8

% ── Mapping ───────────────────────────────────────────────────────────
\subsection{Real-World Mapping}

\begin{longtable}{ll}
\toprule
\textbf{Real-World Concept} & \textbf{Mathematical Object} \\
\midrule
\endhead
problem input & $instance parameters$ \\
\bottomrule
\end{longtable}

% ── Constraints ───────────────────────────────────────────────────────
\subsection{Constraints}

\begin{enumerate}
  \item $Pareto optimality**: A solution is Pareto-optimal if no other feasible solution improves one objective without worsening another$
  \hfill \textit{()}
  \item $Epsilon-constraint method**: Convert a multi-objective problem into a sequence of single-objective problems by constraining all but one objective$
  \hfill \textit{()}
  \item $Knee point**: The point on the Pareto front where the trade-off rate changes most sharply$
  \hfill \textit{(often the best practical compromise)}
  \item $Utopia point**: The ideal (usually infeasible) point formed by the best value of each objective independently$
  \hfill \textit{()}
  \item $Nadir point**: The worst objective values among Pareto-optimal solutions$
  \hfill \textit{(defines the "box" bounding the Pareto front)}
  \item $Marginal rate of substitution**: How much of one objective must be sacrificed to gain a unit of improvement in another$
  \hfill \textit{()}
  \item $Non-dominance**: The defining property of Pareto-optimal solutions$
  \hfill \textit{(no other solution is at least as good in every objective)}
\end{enumerate}

% ── Objective / Claim ─────────────────────────────────────────────────
\subsection{Claim}


% ── Approach ──────────────────────────────────────────────────────────
\subsection{Approach}

\begin{description}
  \item[Suggested approach:] Epsilon-constraint + Pareto filter (scipy L-BFGS-B)
  \item[Complexity class:] 
  \item[Available tools:] Epsilon-constraint + Pareto filter
\end{description}
% ── Phase B: Solution ────────────────────────────────────────────────
\section{Solution}

\begin{polyaresult}

\textbf{Answer}: Solution computed


\textbf{Optimal}: No / Unknown \hfill
\textbf{Feasible}: No
\textbf{Algorithm}: Epsilon-constraint + Pareto filter (scipy L-BFGS-B) \quad $$

\textbf{Time}: 2.315410837996751s


\end{polyaresult}

% ── Solution Details ──────────────────────────────────────────────────
\subsection{Solution Details}


Method: 1. Individual optima: Minimize each objective independently to establish the feasible objective range
2. Epsilon-constraint sweep: Fix one objective as a constraint (durability >= epsilon) and optimize the other (cost), sweeping epsilon across the feasible range
3. Weighted-sum comparison: Solve scalarized problems min w1*f1 + w2*f2 for a grid of weights
4. Pareto filter: Remove dominated points -- a point is dominated if another point is at least as good in all objectives and strictly better in at least one
5. Knee point: Find the point on the Pareto front with maximum curvature (largest angle formed by its neighbors), indicating the best marginal trade-off
6. Utopia/nadir: The utopia point is the (infeasible) vector of individually optimal objectives; the nadir point is the worst objective value attained by any Pareto-optimal solution in each dimension
7. Trade-off rates: Compute the marginal rate of substitution (delta f2 / delta f1) between consecutive Pareto points

% ── Verification ──────────────────────────────────────────────────────
\subsection{Verification}

No independent verification checks recorded.

% ── Phase C: Interpretation ──────────────────────────────────────────
\section{Interpretation}

% ── The Question & The Answer ─────────────────────────────────────────
\subsection{The Question}
A manufacturing company designs a product with two competing objectives: minimize cost and maximize durability (equivalently, minimize inverse durability).

\subsection{The Answer}

\begin{polyainsight}[title={Bottom Line}]
A feasible solution was found.
\end{polyainsight}

% ── What This Means ───────────────────────────────────────────────────
\subsection{What This Means}

- 50 candidate solutions generated via epsilon-constraint sweeps
- Pareto front filtered to \textasciitilde{}20-40 non-dominated points
- Weighted-sum solutions for comparison
- Knee point identified at the point of maximum curvature
- Utopia point (infeasible ideal) and nadir point (worst of the bests)
- Trade-off rates (marginal rate of substitution) along the frontier
- Full independent verification (non-dominance, bounds, objective recomputation)
- Solution saved to solution.json

1. Individual optima: Minimize each objective independently to establish the feasible objective range
2. Epsilon-constraint sweep: Fix one objective as a constraint (durability >= epsilon) and optimize the other (cost), sweeping epsilon across the feasible range
3. Weighted-sum comparison: Solve scalarized problems min w1*f1 + w2*f2 for a grid of weights
4. Pareto filter: Remove dominated points -- a point is dominated if another point is at least as good in all objectives and strictly better in at least one
5. Knee point: Find the point on the Pareto front with maximum curvature (largest angle formed by its neighbors), indicating the best marginal trade-off
6. Utopia/nadir: The utopia point is the (infeasible) vector of individually optimal objectives; the nadir point is the worst objective value attained by any Pareto-optimal solution in each dimension
7. Trade-off rates: Compute the marginal rate of substitution (delta f2 / delta f1) between consecutive Pareto points

% ── Sensitivity Analysis ──────────────────────────────────────────────

% ── Figures ───────────────────────────────────────────────────────────

% ── Recommendations ───────────────────────────────────────────────────
\subsection{Recommendations}

\begin{enumerate}
  \item Review the model assumptions to ensure real-world fidelity.
\end{enumerate}

% ── Limitations ───────────────────────────────────────────────────────
\subsection{Limitations}

\begin{polyawarnbox}
\begin{itemize}
  \item Model assumes deterministic parameters; real-world variability not captured.
  \item Solution quality depends on the accuracy of input data.
\end{itemize}
\end{polyawarnbox}

% ── Appendix ─────────────────────────────────────────────────────────

\end{document}